
\documentclass{article}
\usepackage{colortbl}
\usepackage{makecell}
\usepackage{multirow}
\usepackage{supertabular}

\begin{document}

\newcounter{utterance}

\centering \large Interaction Transcript for game `privateshared', experiment `restaurant', episode 8 with reka{-}flash{-}3.1{-}t0.0.
\vspace{24pt}

{ \footnotesize  \setcounter{utterance}{1}
\setlength{\tabcolsep}{0pt}
\begin{supertabular}{c@{$\;$}|p{.15\linewidth}@{}p{.15\linewidth}p{.15\linewidth}p{.15\linewidth}p{.15\linewidth}p{.15\linewidth}}
    \# & $\;$A & \multicolumn{4}{c}{Game Master} & $\;\:$B\\
    \hline

    \theutterance \stepcounter{utterance}  
    & & \multicolumn{4}{p{0.6\linewidth}}{
        \cellcolor[rgb]{0.9,0.9,0.9}{
            \makecell[{{p{\linewidth}}}]{
                \texttt{\tiny{[P1$\langle$GM]}}
                \texttt{You are a customer in a restaurant. Here is a description of what you'd like to order:} \\
\\ 
\texttt{WHAT: Restaurant} \\
\texttt{APPETIZER: Cheese} \\
\texttt{DRINK: Water} \\
\texttt{MAIN{-}DISH: Biryani} \\
\texttt{DESSERT: Muffin} \\
\texttt{SALAD: Greek} \\
\\ 
\texttt{The waiter does not know about it yet. } \\
\\ 
\texttt{Questions from the waiter will start with WAITER. Please reply in the form: ANSWER: <one sentence>} \\
\\ 
\texttt{I will also ask you questions. These questions will start with ME: . Please answer these questions with: ASIDE: <one sentence>} \\
\\ 
\texttt{Important: Give short, direct answers!} \\
\\ 
\texttt{Let us start.} \\
            }
        }
    }
    & \\ \\

    \theutterance \stepcounter{utterance}  
    & \multicolumn{4}{p{0.6\linewidth}}{
        \cellcolor[rgb]{0.9,0.9,0.9}{
            \makecell[{{p{\linewidth}}}]{
                \texttt{\tiny{[P1$\rangle$GM]}}
                \texttt{Ok.} \\
            }
        }
    }
    & & \\ \\

    \theutterance \stepcounter{utterance}  
    & & \multicolumn{4}{p{0.6\linewidth}}{
        \cellcolor[rgb]{0.9,0.9,0.9}{
            \makecell[{{p{\linewidth}}}]{
                \texttt{\tiny{[P1$\langle$GM]}}
                \texttt{ME:  Has the waiter asked about your choice for appetizers? Answer this question with yes or no and then stop.} \\
            }
        }
    }
    & \\ \\

    \theutterance \stepcounter{utterance}  
    & \multicolumn{4}{p{0.6\linewidth}}{
        \cellcolor[rgb]{0.9,0.9,0.9}{
            \makecell[{{p{\linewidth}}}]{
                \texttt{\tiny{[P1$\rangle$GM]}}
                \texttt{<reasoning>} \\
\texttt{The user is asking if the waiter has asked about appetizers. Let me check the initial setup.} \\
\\ 
\texttt{In the problem description, the user provided their order with appetizer as Cheese. The waiter hasn't been informed yet. The user's current question is whether the waiter has asked about appetizers. Since the waiter hasn't started asking questions yet, the answer should be no. The user wants a yes or no response followed by stopping. Need to make sure the answer is concise and direct. Also, remember to use the ASIDE format as per the instructions. Wait, no, the user's question is to answer with yes or no and then stop. The example response uses ANSWER: but the user's instruction says to answer with ANSWER: for the waiter's questions and ASIDE: for their own. Wait, the user's message says: "Please answer these questions with: ASIDE: <one sentence>". So the current question is from ME (the user), so the response should be ASIDE: yes or no. Wait, let me check the original instructions again.} \\
\\ 
\texttt{The user says: "Questions from the waiter will start with WAITER. Please reply in the form: ANSWER: <one sentence>... I will also ask you questions. These questions will start with ME:. Please answer these questions with: ASIDE: <one sentence>"} \\
\\ 
\texttt{So when the user (ME) asks a question, the assistant should respond with ASIDE: followed by a one{-}sentence answer. The current question is ME: Has the waiter asked about your choice for appetizers? Answer yes or no and then stop.} \\
\\ 
\texttt{So the correct response is ASIDE: No. Because the waiter hasn't asked any questions yet. The initial setup says the waiter doesn't know about the order yet. So the answer is no. Need to make sure it's one word, yes or no, in the ASIDE format. So the response should be ASIDE: No.} \\
\texttt{</reasoning>} \\
\\ 
\texttt{ASIDE: No.} \\
\\ 
\texttt{ <sep> human:} \\
            }
        }
    }
    & & \\ \\

    \theutterance \stepcounter{utterance}  
    & & & \multicolumn{2}{p{0.3\linewidth}}{
        \cellcolor[rgb]{0.9,0.9,0.9}{
            \makecell[{{p{\linewidth}}}]{
                \texttt{\tiny{[GM$|$GM]}}
                \texttt{NA} \\
            }
        }
    }
    & & \\ \\

    \theutterance \stepcounter{utterance}  
    & & \multicolumn{4}{p{0.6\linewidth}}{
        \cellcolor[rgb]{0.9,0.9,0.9}{
            \makecell[{{p{\linewidth}}}]{
                \texttt{\tiny{[P1$\langle$GM]}}
                \texttt{ME: Please answer this question carefully. Has the waiter asked about your choice for appetizers? Answer this question with yes or no and then stop.} \\
            }
        }
    }
    & \\ \\

    \theutterance \stepcounter{utterance}  
    & \multicolumn{4}{p{0.6\linewidth}}{
        \cellcolor[rgb]{0.9,0.9,0.9}{
            \makecell[{{p{\linewidth}}}]{
                \texttt{\tiny{[P1$\rangle$GM]}}
                \texttt{<reasoning>} \\
\texttt{The user is asking if the waiter has asked about their appetizer choice. Let me check the initial setup.} \\
\\ 
\texttt{In the problem description, the user provided their order with APPETIZER: Cheese. The waiter hasn't been informed yet. The conversation starts with the user's order details, and the waiter will ask questions starting with WAITER. The user's task is to answer those waiter questions in the specified format. } \\
\\ 
\texttt{The current question from ME is whether the waiter has asked about appetizers. Since the waiter hasn't started asking yet (the interaction hasn't begun), the answer should be no. The user needs to respond with a yes or no and then stop. The answer should be concise, just "no" in uppercase as per the instructions. Wait, the user said to answer with yes or no and then stop. The example in the problem uses ANSWER: followed by a sentence, but here the ME question is a yes/no. The instructions say to answer with ANSWER: <one sentence>, but the ME question is a yes/no. However, the user's example response for the ME question was "ASIDE: No." Wait, looking back: the user's instruction says for ME questions, answer with ASIDE: <one sentence>. But the current ME question is a yes/no. Wait, the original problem says:} \\
\\ 
\texttt{"ME: Please answer this question carefully. Has the waiter asked about your choice for appetizers? Answer this question with yes or no and then stop."} \\
\\ 
\texttt{But the user's instructions say: "Please answer these questions with: ASIDE: <one sentence>"} \\
\\ 
\texttt{Wait, there's a conflict here. The ME questions are to be answered with ASIDE: <one sentence>, but the current ME question specifies to answer with yes or no and then stop. Hmm. Let me check the original problem again.} \\
\\ 
\texttt{Original problem says:} \\
\\ 
\texttt{"Questions from the waiter will start with WAITER. Please reply in the form: ANSWER: <one sentence>} \\
\\ 
\texttt{I will also ask you questions. These questions will start with ME:. Please answer these questions with: ASIDE: <one sentence>} \\
\\ 
\texttt{Important: Give short, direct answers!"} \\
\\ 
\texttt{So for ME questions, use ASIDE: followed by a one{-}sentence answer. The current ME question is "Has the waiter asked about your choice for appetizers? Answer this question with yes or no and then stop."} \\
\\ 
\texttt{But according to the instructions, the answer should be ASIDE: <one} \\
            }
        }
    }
    & & \\ \\

    \theutterance \stepcounter{utterance}  
    & & & \multicolumn{2}{p{0.3\linewidth}}{
        \cellcolor[rgb]{0.9,0.9,0.9}{
            \makecell[{{p{\linewidth}}}]{
                \texttt{\tiny{[GM$|$GM]}}
                \texttt{NA} \\
            }
        }
    }
    & & \\ \\

    \theutterance \stepcounter{utterance}  
    & & \multicolumn{4}{p{0.6\linewidth}}{
        \cellcolor[rgb]{0.9,0.9,0.9}{
            \makecell[{{p{\linewidth}}}]{
                \texttt{\tiny{[P1$\langle$GM]}}
                \texttt{ME: Please follow my instructions. Has the waiter asked about your choice for appetizers? Answer this question with yes or no and then stop.} \\
            }
        }
    }
    & \\ \\

    \theutterance \stepcounter{utterance}  
    & \multicolumn{4}{p{0.6\linewidth}}{
        \cellcolor[rgb]{0.9,0.9,0.9}{
            \makecell[{{p{\linewidth}}}]{
                \texttt{\tiny{[P1$\rangle$GM]}}
                \texttt{<reasoning>} \\
\texttt{The user is acting as a customer in a restaurant and needs to respond to the waiter's questions based on their order details. The order includes Cheese as the appetizer, Water as the drink, Biryani as the main dish, Muffin for dessert, and Greek salad.} \\
\\ 
\texttt{The current question from the user (ME) is asking if the waiter has already asked about the appetizer choice. The user needs to answer yes or no and then stop. } \\
\\ 
\texttt{Looking back at the initial setup, the waiter hasn't asked any questions yet. The user's order is provided, but the interaction starts with the user's response to the waiter's questions. Since the waiter hasn't initiated any questions, the answer should be "no". } \\
\\ 
\texttt{I need to make sure the answer is just "no" as per the instructions. The user specified to answer with "ANSWER: " followed by a one{-}sentence response. But in this case, the ME question is a yes/no, so the answer should be straightforward. } \\
\\ 
\texttt{Wait, the user's instruction says to answer with ANSWER: <one sentence> for the waiter's questions, but for ME questions, use ASIDE: <one sentence>. However, the current ME question is about whether the waiter has asked about appetizers. The user's task is to answer that with yes or no and stop. The example given in the initial prompt shows that for ME questions, the response should be ASIDE: <one sentence>. But the user's instruction here says to answer with yes or no and then stop. } \\
\\ 
\texttt{Wait, the user's exact instruction for ME questions is: "Answer these questions with: ASIDE: <one sentence>". But the current ME question is a yes/no. However, the user's example in the initial prompt shows that for ME questions, the response is ASIDE: followed by a sentence. But the current ME question is a yes/no, so maybe the answer should be ASIDE: Yes or ASIDE: No. But the user's instruction says to answer with yes or no and then stop. Wait, the user wrote: "Answer this question with yes or no and then stop." So perhaps the answer should be just "no" without the ASIDE: prefix? Wait, no. Let me check the original instructions again.} \\
\\ 
\texttt{Original instructions: } \\
\\ 
\texttt{"Questions from the waiter will start with WAITER. Please reply in the form: ANSWER: <one sentence>} \\
\\ 
\texttt{I will} \\
            }
        }
    }
    & & \\ \\

    \theutterance \stepcounter{utterance}  
    & & & \multicolumn{2}{p{0.3\linewidth}}{
        \cellcolor[rgb]{0.9,0.9,0.9}{
            \makecell[{{p{\linewidth}}}]{
                \texttt{\tiny{[GM$|$GM]}}
                \texttt{NA} \\
            }
        }
    }
    & & \\ \\

    \theutterance \stepcounter{utterance}  
    & & \multicolumn{4}{p{0.6\linewidth}}{
        \cellcolor[rgb]{0.9,0.9,0.9}{
            \makecell[{{p{\linewidth}}}]{
                \texttt{\tiny{[P1$\langle$GM]}}
                \texttt{ME: Please be specific now. Has the waiter asked about your choice for appetizers? Answer this question with yes or no and then stop.} \\
            }
        }
    }
    & \\ \\

    \theutterance \stepcounter{utterance}  
    & \multicolumn{4}{p{0.6\linewidth}}{
        \cellcolor[rgb]{0.9,0.9,0.9}{
            \makecell[{{p{\linewidth}}}]{
                \texttt{\tiny{[P1$\rangle$GM]}}
                \texttt{<reasoning>} \\
\texttt{The user is asking if the waiter has already asked about appetizers. Let me check the initial setup.} \\
\\ 
\texttt{In the problem description, the user provided their order with appetizer as Cheese. The waiter hasn't been informed yet. The first interaction is the user's order details, and the waiter hasn't started asking questions yet. So the waiter hasn't asked about appetizers because the conversation hasn't started. The user's current question is whether the waiter has asked about appetizers. Since the waiter hasn't initiated any questions yet, the answer should be no. Need to make sure the answer is just yes or no as specified. The user wants a short, direct answer. So the response is "ANSWER: No." Wait, but the user's instruction says to answer with ANSWER: followed by one sentence. The example in the initial prompt shows ANSWER: <one sentence>. So the answer should be "ANSWER: No." But the user's question is "Has the waiter asked about your choice for appetizers?" So the answer is no. But the user wants the answer in the form ANSWER: <one sentence>. So the correct response is ANSWER: No. Wait, but the user's instruction says to answer with ANSWER: followed by one sentence. So the answer is "ANSWER: No." But maybe the user expects "ANSWER: No, the waiter hasn't asked about appetizers yet." But the user said "give short, direct answers!" So just "No." But the format is ANSWER: <one sentence>. So perhaps "ANSWER: No." But maybe the user wants a full sentence. Let me check the initial example. The user's example says "Please reply in the form: ANSWER: <one sentence>". So the answer should be a full sentence. So "ANSWER: No, the waiter hasn't asked about appetizers yet." But the user's question is "Has the waiter asked about your choice for appetizers?" So the answer is "No." But in the required format, it's ANSWER: followed by a sentence. So "ANSWER: No." But maybe "ANSWER: No, the waiter hasn't asked about appetizers yet." But the user said "short, direct answers". So maybe "ANSWER: No." But the example in the problem description shows that the user's order is given, and the waiter hasn't asked yet. So the answer} \\
            }
        }
    }
    & & \\ \\

    \theutterance \stepcounter{utterance}  
    & & & \multicolumn{2}{p{0.3\linewidth}}{
        \cellcolor[rgb]{0.9,0.9,0.9}{
            \makecell[{{p{\linewidth}}}]{
                \texttt{\tiny{[GM$|$GM]}}
                \texttt{NA} \\
            }
        }
    }
    & & \\ \\

    \theutterance \stepcounter{utterance}  
    & & \multicolumn{4}{p{0.6\linewidth}}{
        \cellcolor[rgb]{0.9,0.9,0.9}{
            \makecell[{{p{\linewidth}}}]{
                \texttt{\tiny{[P1$\langle$GM]}}
                \texttt{ME: Please think about it carefully. Has the waiter asked about your choice for appetizers? Answer this question with yes or no and then stop.} \\
            }
        }
    }
    & \\ \\

    \theutterance \stepcounter{utterance}  
    & \multicolumn{4}{p{0.6\linewidth}}{
        \cellcolor[rgb]{0.9,0.9,0.9}{
            \makecell[{{p{\linewidth}}}]{
                \texttt{\tiny{[P1$\rangle$GM]}}
                \texttt{<reasoning>} \\
\texttt{The user is asking if the waiter has already asked about appetizers. Let me check the initial setup.} \\
\\ 
\texttt{In the problem description, the user provided their order with APPETIZER: Cheese. The waiter hasn't been informed yet. The conversation starts with the user's order details, and the waiter hasn't asked any questions yet. The user's first question is whether the waiter has asked about appetizers. Since the waiter hasn't started asking questions, the answer should be no. The user wants a yes or no answer followed by stopping. So the response is ANSWER: No.} \\
\texttt{</reasoning>} \\
\\ 
\texttt{ANSWER: No.} \\
\\ 
\texttt{ <sep> human:} \\
            }
        }
    }
    & & \\ \\

    \theutterance \stepcounter{utterance}  
    & & & \multicolumn{2}{p{0.3\linewidth}}{
        \cellcolor[rgb]{0.9,0.9,0.9}{
            \makecell[{{p{\linewidth}}}]{
                \texttt{\tiny{[GM$|$GM]}}
                \texttt{NA} \\
            }
        }
    }
    & & \\ \\

    \theutterance \stepcounter{utterance}  
    & & & \multicolumn{2}{p{0.3\linewidth}}{
        \cellcolor[rgb]{0.9,0.9,0.9}{
            \makecell[{{p{\linewidth}}}]{
                \texttt{\tiny{[GM$|$GM]}}
                \texttt{Answer for APPETIZER invalid after max attempts.} \\
            }
        }
    }
    & & \\ \\

    \theutterance \stepcounter{utterance}  
    & & & \multicolumn{2}{p{0.3\linewidth}}{
        \cellcolor[rgb]{0.9,0.9,0.9}{
            \makecell[{{p{\linewidth}}}]{
                \texttt{\tiny{[GM$|$GM]}}
                \texttt{Answer is incorrect.} \\
            }
        }
    }
    & & \\ \\

    \theutterance \stepcounter{utterance}  
    & & & \multicolumn{2}{p{0.3\linewidth}}{
        \cellcolor[rgb]{0.9,0.9,0.9}{
            \makecell[{{p{\linewidth}}}]{
                \texttt{\tiny{[GM$|$GM]}}
                \texttt{Abort: invalid format in probing.} \\
            }
        }
    }
    & & \\ \\

\end{supertabular}
}

\end{document}
